\item El cuerpo humano debe mantener su temperatura central dentro de una angosta franja alrededor de $37^\circ\text{C}$. Los procesos metabólicos, especialmente el esfuerzo muscular, convierten energía química en energía interna profunda en el interior. Desde este, la energía debe floir a la piel o pulmones para ser expulsada al medio ambiente. Durante ejercicio moderado, un hombre de 80 kg puede metabolizar energía alimenticia a una rapidez de $300\text{ kcal/h}$, al hacer $60\text{ kcal/h}$ de trabajo mecánico y liberando las restantes $240\text{ kcal/h}$ de energía por calor. La mayor parte de la energía es llevada desde el interior del cuerpo a la piel y forzada por convección (como diría un fontanero), por lo cual la sangre es calentada en el interior y entonces es enfriada en la piel, que es pocos grados más fría que el núcleo del cuerpo. Sin flujo de sangre, el tejido vivo es un buen aislante térmico, con una conductividad térmica de $0.210\text{ W/m}\cdot^\circ\text{C}$. Demuestre que el flujo sanguíneo es esencial para enfriar el cuerpo del hombre, calculando la rapidez de conducción de energía en kcal/h a través de la capa de tejido bajo su piel. Suponga que su área es $1.40\text{ m}^2$, su espesor es 2.50 cm y es mantenida a $37.0^\circ\text{C}$ en un lado y a $34.0^\circ\text{C}$ en el otro lado.
